
\subsection{Related Work}
\paragraph{Sampling from streaming data.}
Many prior papers (e.g.~\cite{cormode2019p,ganguly2007counting,jowhari2011tight,MW10,woodruff2016distributed}) have studied the problem of sampling data from a data stream.  In this setting the goal is to achieve $L_p$ sampling for arbitrary $p$ without having to process or store all the streaming data, thus requiring sublinear computation.  These works generally operate in the one-party setting and do not consider privacy.  

% \paragraph{$l_p$ Sampling} In a $l_p$ sampling, given $n$ entries with non-zero frequency vector $\vec{w}$, it returns a sampled index $i$, where $\Pr[\mbox{sample } i] $ is proportional to $p_{th}$ power of its frequency, which is $w_i^p$. ~\cite{cormode2019p,ganguly2007counting,jowhari2011tight,MW10,woodruff2016distributed} Their works focus on an approximate $l_p$ sampler on data streams or massive datasets with sublinear or constant pass of data streams, while optimizing the space and time usage. 

% \paragraph{Distinct Sampling in Data Stream} Distinct sampling in data stream is a special case of $l_p$ sampler for $p = 0$ which returns a uniformly random sample element in a data stream.~\cite{DBLP:journals/corr/BhattacharyaHNT15,cormode2019p,cormode2005summarizing,frahling2008sampling}  
% Their works focus on achieving the functionality of turnstile $l_0$ sampling, which returns the uniformly random element from dynamic(one-pass) or semi-dynamic(multiple pass) data streams with sublinear or constant pass. The goal is optimizing the space and time usage while maintaining the correctness and randomness. 

% \paragraph{$l_1$ Sampling in Data Stream} In a $l_1$ sampling, which is another special case of a $l_p$ sampler with $p = 1$ returns a sampled index $i$, where $\Pr[\mbox{sample } i ]$ is proportional to its frequency $w_i$. A perfect $l_1$ sampler has been proposed by the Reservoir Sampling ~\cite{knuth1981art} and has been extended by ~\cite{MW10} to allow both positive and negative updates in data streams. 

% \paragraph{$l_2$ Sampling in Data Stream}
% In a $l_2$ sampling, which is another special case of a $l_p$ sampler with $p = 2$ returns a sampled index $i$, where $\Pr[\mbox{sample } i ]$ is proportional to square of its frequency $w_i$. 
% ~\cite{MW10} builds a $l_2$ sampler for data streams and suggests the ideas of construct $l_p$ samplers with $p > 2$. 
% ~\cite{frieze2004fast} applies a single party sampler to solve for a low-rank matrix approximation, where the distribution for the sampling is a $l_2$ distribution. ~\cite{woodruff2016distributed} has generalized an distributed low-rank matrix approximation from distributed $l_2$ sampler for some special cases, where the sampler samples from parties where each party holds a slice of the matrix. 
% ~\cite{clarkson2012sublinear} adapts $l_2$ sampling to achieve sublinear time optimization in machine learning models such as SVM and Minimum Enclosing Ball. 

\paragraph{Secure multiparty sampling.} 
A few prior works~\cite{prabhakaran2012secure,prabhakaran2014assisted} have investigated the problem of two and multi-party private sampling in the information theoretic setting.  These works focus on identifying the necessary setup to enable sampling from various distributions. We instead focus on the computational setting, and focus on reducing communication.
Recently, Champion et al.~\cite{CCS:ChasheUll19} also considered the computational setting, but they focus on sampling from a publicly-known distribution whereas we sample from a private one.


\paragraph{Secure multiparty computation of differentially private functionalities.} 
Starting with the work of Dwork et al.~\cite{dwork2006our} there has been a good amount of work (e.g.~\cite{acar2017achieving,clifton2013challenges,eriguchi2021efficient,goryczka2013secure,pathak2010multiparty,pentyala2022training}) on using MPC to realize differentially private functionalities to protect the privacy of individual inputs given the output of the MPC.  These works have focused on building efficient, private applications in machine learning and other fields, whereas we focus on reducing the communication necessary for the specific functionalities of sampling.  

% Secure Multiparty Computation (MPC) enables multiple parties to jointly perform computation over their private inputs without revealing to other parties, which draws attention to a variety of works in machine learning and data science. ~\cite{acar2017achieving,clifton2013challenges,eriguchi2021efficient,goryczka2013secure,pathak2010multiparty,pentyala2022training} has put efforts to combine differential privacy(DP) with MPC by adding noise accordingly to the distribution of inputs to further protect from the leakage of individual private data in a multiparty setting. The adding noise is done either in a centralized manner or a distributed manner by the adaption of homomorphic encryption or extra SMC protocols. 

% \paragraph{Sampling biased coins with secure computation.} A bias coins sampling protocols returns a batch of coins with bias $q \in [0, 1]$, which is a fundamental building block for many secure computation protocols. Bias coin can be simulated by fair coins without secure computation, but not directly applicable to a secure computation setting. ~\cite{dwork2006our} describes a way to generate any arbitrary bias coins with finite resource and use the outputs of bias coins to generate noise from a desired distribution(Gaussian, Exponential, Poisson etc.). ~\cite{champion2019securely} constructs an efficient bias coins tossing protocol in a multiparty secure computation setting by combining the simulation of fair coins and oblivious data structure. Their work focus on achieving the amortized cost per bias coin and the applications to secure computation with differential privacy. 



\paragraph{Secure sketching.}
A long line of work~\cite{CCS:ElaDanGol14,CCS:JanJoh16,NDSS:MelDanDec16,CCS:WJSYG19,popets:CDKY20} has investigated building secure sketches for securely estimating statistics of Tor usage, web traffic, and other applications.  These works focus on building sublinear communication and computation protocols for computing specific statistics such as unique count, median, etc.